\section{Limits, open problems and status}\label{sec:limits}
\sectionbrief
{all earlier sections.}
{to separate in one list what the theory forbids, what is open in the theory, and what is merely not finished in the code.}
{honest summary for a reader who needs to decide what can be trusted.}

\subsection{What the theory forbids}
\begin{itemize}[nosep]
\item A fit from finite data cannot be certified; the output is a label and a posterior, never ``identified'' \thy{T17}.
\item No data-free noise threshold exists \thy{T15}; thresholds are posterior quantities.
\item No method beats the error floor $\gamma_N^{-1/2}\sigma$ on components \thy{T6}; the expected worst-case noise rate for coefficients is $\varepsilon^{1/3}$ \thy{T14}.
\item $\Eint$ and the annihilator are discontinuous: uncertainty only by Monte Carlo \thy{T12}.
\item The factor split fails on domains with holes: rectangles only \thy{T1}.
\item The factor split loses one derivative \thy{T11}: component priors must be one order smoother than the field.
\end{itemize}

\subsection{What is open in the theory}
Conditioning of the Hankel/Prony step under noise (\thy{T10}, open): the practical noise threshold is empirical. A repeated factor has no finite test (it is reported \textsc{undecided}). The closure question for $(\tau-\xi^2)(\tau+c\xi)$ is open. The algorithm covers only constant coefficients in $1+1$ dimensions; variable coefficients, nonlinear equations (Burgers, KdV), fractional and Dirac operators and higher dimensions are outside it.

\subsection{What is not finished in the code}
\begin{table}[!htb]
\centering\small
\begin{tabular}{@{}L{4.6cm}L{2.4cm}L{5.8cm}@{}}
\toprule
Item & Status & Note\\\midrule
Core v1, core v2 and its extensions (per-bin noise, higher orders, confluent test) & \impl{} & committed; full suite reported as $238$ passed at the last handoff\\
Aggregation, DP, clustering, domain, baselines & \impl{} & committed in Phase~2\\
GP interpolation & \impl{} & band coverage within $\pm5$ points in $11$ of $12$ cells\\
Friedrichs layer (\texttt{bayes\_split}) & \impl{} & logged coverage $0.85$--$1.00$; criterion not certified\\
Loop and refinement (Sections~\ref{sec:loop}, \ref{sec:refine}) & \impl{} & committed; thresholds hand defaults ($P_{\mathrm{conc}}$, floor $0.02$, degree limits, $1$ nat, $-3$ in $\log_{10}p$)\\
Scattered data end to end & \impl{} & final benchmark $30/30$ right (below)\\
Held-out random equations & \impl{} (measured) & $20$ draws: $12$ right, $7$ equivalent, $1$ abstention, no false-confident\\
Learned layers & evaluated & none ships; an exact memo (no change of output) is on; see the companion document\\
Controller, L1b, L6, calibration of the hand defaults & \plan{} & not built\\
\bottomrule
\end{tabular}
\caption{Status of each component at the time of writing.}
\end{table}

\subsection{The benchmarks, from the first run to the final one}
\textbf{Run 6 (before the later fixes).} Ten named cases, $30$ runs per cell: right factor set at noise $10^{-3},\,10^{-2},\,5\cdot10^{-2}$ on the grid $1.00,\,1.00,\,0.93$; at $2000$ scattered points $1.00,\,0.97,\,0.90$; at $400$ points $1.00,\,0.93,\,0.73$; at least as good as every sparse-regression baseline in all nine cells; median run time about $5$ s.

\textbf{Final benchmark (P8).} Ten cases $\times$ \{grid $1\%$, grid $5\%$, $400$ points $5\%$\}, one seed: the pipeline gives the right factor set on $\mathbf{30/30}$ inputs, no false-confident answer (run 6 on the same inputs: $26$ right, $3$ abstentions, $1$ false-confident). The sparse-regression baselines on the same inputs: weak SINDy $23$ right and $5$ false-confident, ensemble SINDy $19$ and $8$, Bayesian SINDy $20$ and $10$, PDE-FIND $4$ and $25$, BOP-DMD $1$ and $29$. The coverage of the coefficient intervals is $0.94$ at nominal $0.90$ on all $30$ inputs; the weight intervals (only $6$ inputs have two factors) cover $1.00$. The price: the median run takes about $127$ s and the total about $1.9$ hours, against $5$ s and $26$ minutes in run 6, and single scattered cases at $400$ points take tens of minutes.

\textbf{Equations nobody tuned for (E13).} Random linear equations from five families (the simulator of Section~\ref{sec:train}), $20$ draws at noise $10^{-2}$: $12$ right, $7$ equivalent (the data cannot tell them from the truth), $1$ correct abstention, \emph{no false-confident answer}; before the later fixes: $10$ right, $9$ equivalent, $1$ abstention and $1$ false-confident. The sparse-regression baselines get at most $6$ right. On the hand-picked held-out suites: $13$ of $14$ right and $11$ of $11$ on conjugate pairs. These are small samples: indications, not rates.

\subsection{Three cautions on the evidence}
First, the ten named cases tuned the defaults and are also the cases that measure them; the random draws above are held-out evidence but few, and the scattered seeds of the benchmark share one simulated field, so repeated misses are not independent evidence. Second, grouping results use exact components and boundaries aligned with windows. Third, the toy example of the first note ($0.3e^{-t}\cos x$) has a single mode for its diffusion factor, for which the coefficient is not identifiable; the examples here use many modes.

\begin{keypoint}
The algorithm is solid where its inputs are gridded and its structure is linear, constant-coefficient and $1+1$. The route from scattered noisy data through the GP posterior to the refined, honestly labelled answer works on the named cases and on the first random draws, at a cost in time. The open steps are generality on larger samples, two conjugate pairs (found but labelled \textsc{undecided}), calibrated weight intervals over the ambiguity of a shared mode, replacing the hand defaults by inferred values, and speed.
\end{keypoint}
